\section{The cost at the first-exit board}\label{costes}

Let $B$ be the first-exit board of a history, and let $k$ be its size. When $Q(B)>0$, recall that
$g(B)=\log Q(B)-L(B)$ and set $g^+(B)=\max\{g(B),0\}$. If $Q(B)=0$, set $g^+(B)=0$.
We call $g^+(B)$ the \emph{exit cost}. All estimates in this section concern a fixed board.
The diffusion, vacancy, and gap bounds we use are the outer bounds that the board inherits from its
regular parent. We do not assume the regularity conditions that the board has just failed to satisfy.

\subsection{The normalizer on a residual board}\label{costes:normalizador}

Let $p=(p_e)$ be the Bethe point of $B$. Each square $e=(i,j)$
lies on a row, a column, and two diagonals $D(e),A(e)$. As in Section~\ref{leyes},
$\ell_f=\sum_{e\in f}p_e$ and $h_f=1-\ell_f$ are the mass and vacancy of each residual diagonal $f$.
Write $\mathcal A=(A_R,D)$ for the incidence matrix with squares indexing the rows and lines indexing
the columns. Its first columns correspond to the board rows and columns, and the remaining columns
to the diagonals. Thus $D$ is the square--diagonal incidence matrix of
Section~\ref{macro:reconstruccion}, and $A_R=B^T$ is the transpose of the row--column incidence
matrix in that section. Here we reserve $B$ for the board. Define
\[
 \tau_d(p)=\frac{p(1-p)}{1-dp},\qquad
 C_d=\sum_e\bigl[\log(1-dp_e)-d\log(1-p_e)\bigr],\qquad d=2,4 .
\]
In the asymptotic estimates we take $k$ large enough that $4p_{\max}\le1/2$.
For $d=4$ the matrix in the normalizer is
\[
 S_4=\operatorname{diag}(0_{RC},h)+
          \mathcal A^T\operatorname{diag}(\tau_4(p_e))\mathcal A .
\]
This is the line operator $S_B$ of Section~\ref{leyes}. Its zero mode is
$(\mathbf1_R,-\mathbf1_C,0)$. The pseudodeterminant $\pdet$ is the product of the positive
eigenvalues. In this notation the predictor takes the form
\[
 \Gamma_4=\frac{e^{-C_4/2}\sqrt{4\pi k}}{\sqrt{\pdet S_4}},
 \qquad L(B)=H_B(p)+\log\Gamma_4 ,
\]
where $H_B$ is the Bethe entropy \eqref{bethe}. The hypotheses of the following lemma are the
diffusion, vacancy, and gap bounds from Section~\ref{leyes}: the matrix
$A_R^T\operatorname{diag}(p)A_R$ is the row--column Laplacian up to a sign change in the
column coordinates.

\begin{lema}[Estimate for the normalizer]\label{coste-estatico}
Suppose that $p_{\max}\le C/k$, that $h_f\ge\eta>0$, and that
$A_R^T\operatorname{diag}(p)A_R$ has a gap of at least $\gamma>0$ on the orthogonal complement
of $(\mathbf1_R,-\mathbf1_C)$. Then
\[
 L(B)-H_B(p)=\tfrac12\log(4\pi k)+O_{C,\eta,\gamma}(1).
\]
\end{lema}
\begin{proof}
Write $\tau_e=\tau_4(p_e)$ and $E=S_4-I$. The sum of the incident $p_e$ is one on a row or column
and $1-h_f$ on a diagonal. Hence, for every line $f$,
\[
 E_{ff}=\sum_{e\in f}(\tau_e-p_e),\qquad
 \tau_e-p_e=\frac{3p_e^2}{1-4p_e}.
\]
Since $\sum_e p_e=k$ and $\sum_e p_e^2\le C$, we obtain
$\operatorname{tr}E=4\sum_e(\tau_e-p_e)=O(1)$. Moreover, $0\le E_{ff}=O(1/k)$ because the mass of
each line is at most one. Consequently $\sum_f E_{ff}^2=O(1/k)$. Two lines from different
families intersect in at most one square. Each square gives twelve ordered pairs of distinct lines, so
\[
 \sum_{f\ne j}E_{fj}^2=12\sum_e\tau_e^2=O(1).
\]

Let $y\perp(\mathbf1_R,-\mathbf1_C)$, let $z$ be a vector indexed by diagonals, and set
$P_0=\operatorname{diag}(p)$. Since $\tau_e\ge p_e$, the quadratic form of $S_4$ at $(y,z)$
dominates $\|P_0^{1/2}(A_Ry+Dz)\|^2+\eta\|z\|^2$. Each square lies on two diagonals, so
\[
 \|P_0^{1/2}Dz\|^2
 \le2\sum_f\ell_f z_f^2\le2\|z\|^2.
\]
For $0<\theta<1$, completing the square gives
$\|a+b\|^2\ge\theta\|a\|^2-\theta\|b\|^2/(1-\theta)$. With $\theta=\eta/(\eta+4)$, the preceding
quadratic form is at least $\theta\gamma\|y\|^2+(\eta/2)\|z\|^2$.
This proves that the indicated mode is the only zero mode and that the positive eigenvalues are at least
$c=\min\{\theta\gamma,\eta/2\}$. For these eigenvalues we have
\[
 \sum_{\lambda>0}(\lambda-1)=\operatorname{tr}E+1,\qquad
 \sum_{\lambda>0}(\lambda-1)^2\le\operatorname{tr}E^2=O(1).
\]
The Taylor bound $|\log\lambda-(\lambda-1)|\le
(\lambda-1)^2/(2\min\{1,c\}^2)$ gives $\log\pdet S_4=O(1)$.
Finally, the linear terms in $C_4$ cancel, so $|C_4|\le C'\sum_e p_e^2=O(1)$.
Substituting these two estimates into $\log\Gamma_4$ proves the formula.
\end{proof}
